\documentclass[10pt]{extarticle}
\usepackage[a4paper,margin=1.5cm]{geometry}
\usepackage[T1]{fontenc}
\usepackage{amsmath,amssymb}
\usepackage{graphicx}
\usepackage[export]{adjustbox}
\usepackage{multicol}
\usepackage{enumitem}
\usepackage{booktabs,array,tabularx}
\setlength{\columnsep}{0.7cm}
\setlength{\parindent}{0pt}
\setlength{\parskip}{3pt}
\pagestyle{empty}
\makeatletter
% figure registry (figdefs.tex is generated by the extract step)
\newcommand{\qfigdef}[4]{\@namedef{qf@file@#1}{#2}\@namedef{qf@w@#1}{#3}\@namedef{qf@wide@#1}{#4}}
\newcommand{\qf@check}[1]{\@ifundefined{qf@file@#1}{\PackageError{score2}{Unknown figure id #1}{}}{}}
\newcommand{\qf@img}[4][]{\includegraphics[width=#2\dimexpr\@nameuse{qf@w@#4}pt\relax,max width=#3,max height=0.45\textheight,#1]{\@nameuse{qf@file@#4}}}
% Full-width questions: finalize wraps every question that contains a wide figure in \qfullbegin...\qfullend,
% so the WHOLE question (stem, figures, options) is set across the page in reading order, then the two columns
% resume. The question's head - its number, stem and first figure (\qfullsplit marks the end of the head) -
% is kept on one page, so a number is never left at a page foot with its figure overleaf; the rest may break
% across pages normally, so no big gaps are left. Without a \qfullsplit, the whole question is kept together
% if it is under 45% of a page. Never write these yourself.
\newif\ifqf@full
\newif\ifqf@boxed
\newsavebox{\qf@box}
\newcommand{\qf@open}{\global\qf@fulltrue\par\global\qf@boxedtrue
  \setbox\qf@box\vbox\bgroup\hsize\textwidth\linewidth\textwidth}
\newcommand{\qf@emit}[1]{\par\egroup\global\qf@boxedfalse
  \ifdim\dimexpr\ht\qf@box+\dp\qf@box\relax<#1\textheight\noindent\usebox{\qf@box}\par\else\unvbox\qf@box\fi}
\newcommand{\qfullsplit}{\ifqf@boxed\qf@emit{0.85}\fi}
\newcommand{\qf@close}{\ifqf@boxed\qf@emit{0.45}\fi\global\qf@fullfalse}
\newcommand{\qfullbegin}{\end{multicols}\qf@open}
\newcommand{\qfullend}{\qf@close\begin{multicols}{2}}
\newcommand{\qfullstart}{\qf@open}
\newcommand{\qfulldone}{\qf@close}
% \qfig[scale]{id}: block figure, centred, directly where it is written (size set by the tool).
% Fallback if a wide figure is not inside a full-width question: break out of the columns in place - no floats.
\newcommand{\qfig}[2][1]{\qf@check{#2}%
  \ifqf@full\qf@blk{#1}{#2}\else
  \ifnum\@nameuse{qf@wide@#2}=1 \expandafter\@firstoftwo\else\expandafter\@secondoftwo\fi
  {\end{multicols}\begin{center}\qf@img{#1}{\textwidth}{#2}\end{center}\begin{multicols}{2}}%
  {\qf@blk{#1}{#2}}\fi}
\newcommand{\qf@blk}[2]{\par\smallskip{\centering\qf@img{#1}{\linewidth}{#2}\par}\smallskip}
% \qfigi[scale]{id}: inline figure (inside option lists / tables / minipages); never breaks columns.
\newcommand{\qfigi}[2][1]{\qf@check{#2}\qf@img[valign=c]{#1}{\linewidth}{#2}}
\makeatother
% structure macros - lint reads these to check completeness and figure placement
\newcommand{\papersec}[1]{\par\bigskip{\centering\Large\bfseries PAPER-#1\par}\medskip}
\newcommand{\partsec}[2]{\par\medskip\noindent{\large\bfseries PART-#1: #2}\par\smallskip}
\newcommand{\sect}[1]{\par\smallskip\noindent\textbf{#1}\par}
\newcommand{\qn}[1]{\par\smallskip\noindent\textbf{#1.}~}
\newcommand{\paraq}[2]{\par\smallskip\noindent\textbf{Paragraph for Questions #1 and #2}\par}
\newenvironment{qoptions}{\begin{itemize}[leftmargin=1.8em,itemsep=1pt,topsep=2pt,label={}]}{\end{itemize}}
\newcommand{\opt}[2]{\item[(#1)] #2}
% figure registry (generated by score2_tools.py finalize)
\qfigdef{p1_phy_s0_q01_2}{SRGBT - Test 05 (23 Aug 2021)-img/p1_phy_s0_q01_2.png}{51.5}{0}
\qfigdef{p1_phy_s0_q02_1}{SRGBT - Test 05 (23 Aug 2021)-img/p1_phy_s0_q02_1.png}{61.2}{0}
\qfigdef{p1_phy_s0_q03_2}{SRGBT - Test 05 (23 Aug 2021)-img/p1_phy_s0_q03_2.png}{174.6}{0}
\qfigdef{p1_phy_s0_q07_1}{SRGBT - Test 05 (23 Aug 2021)-img/p1_phy_s0_q07_1.png}{126.9}{0}
\qfigdef{p1_phy_s0_q09_1}{SRGBT - Test 05 (23 Aug 2021)-img/p1_phy_s0_q09_1.png}{156.7}{0}
\qfigdef{p1_phy_s0_q10_1}{SRGBT - Test 05 (23 Aug 2021)-img/p1_phy_s0_q10_1.png}{89.3}{0}
\qfigdef{p1_phy_s0_q13_1}{SRGBT - Test 05 (23 Aug 2021)-img/p1_phy_s0_q13_1.png}{196.1}{0}
\qfigdef{p1_phy_s0_q14_1}{SRGBT - Test 05 (23 Aug 2021)-img/p1_phy_s0_q14_1.png}{153.2}{0}
\qfigdef{p1_phy_s0_q16_1}{SRGBT - Test 05 (23 Aug 2021)-img/p1_phy_s0_q16_1.png}{51.0}{0}
\qfigdef{p1_chem_s0_q32_1}{SRGBT - Test 05 (23 Aug 2021)-img/p1_chem_s0_q32_1.png}{86.8}{0}
\raggedcolumns
\makeatletter\@beginparpenalty=10000\makeatother
\newcolumntype{L}[1]{>{\raggedright\arraybackslash}p{#1}}
\begin{document}
\begin{center}{\Large\bfseries SRGBT --- Test 05 (23 Aug 2021)}\end{center}
\begin{multicols}{2}
\raggedcolumns
\papersec{1}
\partsec{1}{PHYSICS}
\qn{1} In the system shown in figure, Both the blocks A and B have equal mass of m and the force constant of the ideal spring is k. Pulley and threads are massless.\par
The time period of vertical oscillations of block A is :\par
\qfig{p1_phy_s0_q01_2}
\begin{qoptions}
\opt{A}{$2\pi\sqrt{\dfrac{3\mathrm m}{7\mathrm K}}$}
\opt{B}{$2\pi\sqrt{\dfrac{4\mathrm m}{3\mathrm K}}$}
\opt{C}{$2\pi\sqrt{\dfrac{5\mathrm m}{4\mathrm K}}$}
\opt{D}{$2\pi\sqrt{\dfrac{5\mathrm m}{2\mathrm K}}$}
\end{qoptions}
\qn{2} In the system shown in the figure the string, springs and pulley are light. The force constant of the two springs are k$_1$ = k and k$_2$ = 2k. Block of mass M is pulled vertically down from its equilibrium position and released. The top surface of the block (represented by line AB) always remains horizontal. The angular frequency of oscillation is found to be $\sqrt{\dfrac{\mathrm{CK}}{3\mathrm m}}$. The value of C is\par
\qfig{p1_phy_s0_q02_1}
\begin{qoptions}
\opt{A}{3}
\opt{B}{5}
\opt{C}{6}
\opt{D}{8}
\end{qoptions}
\qn{3} A wooden platform can be rotated about its vertical axis with constant angular speed $\omega$ with the help of a motor. A buzzer is fixed at the circumference of the platform and it rotates in a circle of radius R. The buzzer produces sound of frequency f$_0$. A mic is placed just beneath the platform near its circumference. An electronic frequency analyzer connected to the mic records the frequency (f) received by the mic. Take time (t) to be zero when the buzzer is just above the mic. Speed of sound = V$_0$. The expression of f as a function of time is found to be\par
$\mathrm f_0\left[\dfrac{\mathrm V_0}{\mathrm V_0+\mathrm R\omega\left(\dfrac{\omega\mathrm t}{\mathrm K}\right)}\right]$. The value of K is\par
\qfig{p1_phy_s0_q03_2}
\begin{qoptions}
\opt{A}{0}
\opt{B}{1}
\opt{C}{2}
\opt{D}{5}
\end{qoptions}
\qn{4} A particle executing SHM with time period T and amplitude A along horizontal direction along x axis. When it is at +$\dfrac{\mathrm A}2$ from equilibrium position and moving towards right, it received an impulse which doubles its velocity without changing its direction.\par
The time taken by particle to reach extreme right position is $\dfrac{\mathrm T}{2\pi}\cos^{-1}\left[\dfrac1{\mathrm K}\right]$. The value of K is
\begin{qoptions}
\opt{A}{8}
\opt{B}{2}
\opt{C}{13}
\opt{D}{3}
\end{qoptions}
\qn{5} Density of a liquid varies with depth as $\rho$ = $\alpha$h. A small ball of density $\rho_0$ is released from the free surface of the liquid. Then the ball will execute SHM of amplitude k times $\dfrac{\rho_0}{\alpha}$. The vaule of k is
\begin{qoptions}
\opt{A}{1}
\opt{B}{2}
\opt{C}{3}
\opt{D}{4}
\end{qoptions}
\qn{6} A toy train in a children amusement park runs on an elliptical orbit having major and minor axis in the ratio of 4 : 3. The length of the train is exactly equal to half the perimeter of the elliptical track. A train is travelling at a constant speed of 20 m/s. The engine sound a whistle when its acceleration is minimum. The frequency of whistle recieved by a passenger sitting in the central compartment of the train is found to be \underline{\hspace{1cm}} (in Hz)\par
(Take speed of sound = 330 m/s; original frequency of whistle = 3460 Hz)
\begin{qoptions}
\opt{A}{3420 Hz}
\opt{B}{3480 Hz}
\opt{C}{3400 Hz}
\opt{D}{3450 Hz}
\end{qoptions}
\qn{7} A uniform rod AB of mass m and length L is tied, at its end B, to a thread which is attached to point P on the ceiling. Length of the thread PB is 0.75 L. The other end A of the rod is hinged at a point on the ceiling. Distance AP = 1.25 L. End B of the rod is pushed gently perpendicular to the plane of the figure and it starts oscillating.\par
\qfig{p1_phy_s0_q07_1}
\begin{qoptions}
\opt{A}{The moment of inertia of the rod about line AP is $\dfrac{13}{25}\mathrm{mL}^2$}
\opt{B}{Assuming that the triangle APB makes a small angle $\theta$ with the vertical plane. The time period of oscillation is $2\pi\sqrt{\dfrac{2\mathrm L}{3\mathrm g}}$}
\opt{C}{Assuming that the triangle APB makes a small angle $\theta$ with the vertical plane. The time period of oscillation is $2\pi\sqrt{\dfrac{2\mathrm L}{5\mathrm g}}$}
\opt{D}{The moment of inertia of rod about line AP is $\dfrac{3\mathrm{mL}^2}{25}$}
\end{qoptions}
\qn{8} The speed v of a particle moving along a straight line, when it is at a distance x from a fixed point on the line is v$^2$ = 144 - 9x$^2$. Select the correct alternative(s) :
\begin{qoptions}
\opt{A}{The motion of the particle is SHM with time period T = $\dfrac{2\pi}3$ unit.}
\opt{B}{The maximum displacement of the particle from the fixed point is 4 unit}
\opt{C}{The magnitude of acceleration at a distance 3 units from the mean position is 27 unit}
\opt{D}{The motion of the particle is periodic but not simple harmonic}
\end{qoptions}
\qn{9} Initially each spring was at natural length. The system is now brought into SHM. So that there is no relative motion between Aand B as well as B and C. Analyse the situation when oscillation is taking place with maximum possible amplitude.\par
\qfig{p1_phy_s0_q09_1}
Choose the correct statement(s)
\begin{qoptions}
\opt{A}{The maximum possible amplitude of oscillation is 2m.}
\opt{B}{The maximum possible friction between A and B during oscillation is 100 N}
\opt{C}{The maximum possible friction between B and C during oscillation is 100 N}
\opt{D}{The time period of oscillation is $\left(\dfrac{2\pi}{\sqrt5}\right)$sec.}
\end{qoptions}
\qn{10} Consider all springs and strings to be ideal. The block is in equilibrium.\par
\qfig{p1_phy_s0_q10_1}
Now the block is displaced downwards by a small amount and released.\par
Choose the correct statement(s)
\begin{qoptions}
\opt{A}{The mechanical energy of oscillation is conserved}
\opt{B}{The time period of oscillation is $2\pi\sqrt{\dfrac{5\mathrm M}{3\mathrm K}}$}
\opt{C}{The time period of oscillation is $2\pi\sqrt{\dfrac{5\mathrm M}{18\mathrm K}}$}
\opt{D}{The frequency at kinetic energy function K(t) is twice the linear frequency at oscillating block.}
\end{qoptions}
\qn{11} The air column in a pipe closed at one end is made to vibrate in its second overtone by a tuning fork of frequency 440 Hz. The speed of sound in air is 330 m/s. P$_0$ is atmospheric pressure in the pipe and $\Delta$P$_0$ is maximum amplitude of pressure variation. Neglect end correction.
\begin{qoptions}
\opt{A}{The length L of air column is (93 $\pm$ 1) cm}
\opt{B}{The amplitude of pressure variation at the middle of the column is $\left[\dfrac{\Delta\mathrm P_0}{\sqrt2}\right]$}
\opt{C}{The length L of air column is (51 $\pm$ 1) cm}
\opt{D}{The amplitude of pressure variation at middle of column is $\left(\dfrac{\Delta\mathrm P}2\right)$}
\end{qoptions}
\qn{12} A source of sound is located in a medium in which speed of sound is V and an observer is located is a medium in which speed of sound is 2V. Both the source and observer are moving directly towards each other at velocity V/5. The source has a frequency of f$_0$.\par
Choose the correct statement(s) :
\begin{qoptions}
\opt{A}{The frequency received by observer is found to be $\dfrac{11}8$f$_0$}
\opt{B}{The frequency received by observer is found to be $\dfrac{13}8$f$_0$}
\opt{C}{The wavelength of sound recieved by observer is $\left(\dfrac{8\mathrm V}{5\mathrm f_0}\right)$}
\opt{D}{The wavelength of sound recieved by observer is $\left(\dfrac{7\mathrm V}{3\mathrm f_0}\right)$}
\end{qoptions}
\qn{13} A wire having mass per unit length $\mu$ and length L is fixed between two fixed vertical walls at a separation L. Due to its own weight the wire sags. The sag in the middle is d ($<<$ L). Assume that tension is practically constant along the wire, owing to its small mass.\par
\qfig{p1_phy_s0_q13_1}
Choose the correct statement(s) :
\begin{qoptions}
\opt{A}{The speed of transverse wave of the wire (approximately) is L$\sqrt{\left(\dfrac{\mathrm g}{8\mathrm d}\right)}$}
\opt{B}{The speed of transverse wave of the wire (approximately) is L$\sqrt{\left(\dfrac{\mathrm g}{4\mathrm d}\right)}$}
\opt{C}{The tension in the string is T = $\left(\dfrac{\mu\mathrm L^2}{8\mathrm d}\right)$g}
\opt{D}{The tension in the string is T = $\left(\dfrac{\mu\mathrm L^2}{4\mathrm d}\right)$g}
\end{qoptions}
\qn{14} There are three sources of sound of equal amplitude with frequency (n $-$ 1), n, (n + 1) Hz. At t = 0, they are all switched on simultaneously.\par
Choose the correct statement(s) :
\begin{qoptions}
\opt{A}{At t = $\dfrac13$ sec, minima will occur}
\opt{B}{At t = $\dfrac23$ sec, minima will occur}
\opt{C}{At t = $\dfrac12$ sec secondary maxima will occur}
\opt{D}{The maximum possible intensity will be 9A$^2$.}
\end{qoptions}
\paraq{15}{16}Initially the system was at rest and all three springs were at their natural lengths. Now each spring is extended by maximum possible x$_0$ and released so that there is no relative motion among A, B, C and D, E, F systems. The entire system oscillates like a two block system.\par
\qfig{p1_phy_s0_q14_1}
$\mu_\mathrm{AB}$ = $\mu_\mathrm{DE}$ = 1.0; $\mu_\mathrm{BC}$ = $\mu_\mathrm{EF}$ = 0.8
\qn{15} The value of x$_0$ is
\begin{qoptions}
\opt{A}{$\dfrac37$m}
\opt{B}{$\dfrac27$m}
\opt{C}{1/7 m}
\opt{D}{1/2 m}
\end{qoptions}
\qn{16} The maximum friction acting B and C during oscillation is
\begin{qoptions}
\opt{A}{(300 $\pm$ 5) N}
\opt{B}{(340 $\pm$ 5) N}
\opt{C}{(280 $\pm$ 5) N}
\opt{D}{(240 $\pm$ 5) N}
\end{qoptions}
\paraq{17}{18}A spring block pendulum is shown in the figure. The system is hanging in equilibrium. A bullet of mass (m/2) moving at a speed u hits the block from downward direction and gets embedded into it.\par
\qfig{p1_phy_s0_q16_1}
\qn{17} The amplitude of oscillation of block is
\begin{qoptions}
\opt{A}{$\sqrt{\dfrac{\mathrm{mu}^2}{2\mathrm k}+\left(\dfrac{\mathrm{mg}}{2\mathrm k}\right)^2}$}
\opt{B}{$\sqrt{\dfrac{\mathrm{mu}^2}{4\mathrm k}+\left(\dfrac{\mathrm{mg}}{2\mathrm k}\right)^2}$}
\opt{C}{$\sqrt{\dfrac{\mathrm{mu}^2}{6\mathrm k}+\left(\dfrac{\mathrm{mg}}{2\mathrm k}\right)^2}$}
\opt{D}{$\sqrt{\dfrac{\mathrm{mu}^2}{8\mathrm k}+\left(\dfrac{\mathrm{mg}}{2\mathrm k}\right)^2}$}
\end{qoptions}
\qn{18} The time taken by block to reach its upper extreme position after hit by bullet is ( A = amplitude)
\begin{qoptions}
\opt{A}{$\sqrt{\dfrac{3\mathrm m}{5\mathrm k}}\cos^{-1}\left(\dfrac{\mathrm{mg}}{3\mathrm{kA}}\right)$}
\opt{B}{$\sqrt{\dfrac{3\mathrm m}{\mathrm k}}\cos^{-1}\left(\dfrac{\mathrm{mg}}{3\mathrm{kA}}\right)$}
\opt{C}{$\sqrt{\dfrac{3\mathrm m}{2\mathrm k}}\cos^{-1}\left(\dfrac{\mathrm{mg}}{2\mathrm{kA}}\right)$}
\opt{D}{$\sqrt{\dfrac{\mathrm m}{2\mathrm k}}\cos^{-1}\left(\dfrac{\mathrm{mg}}{3\mathrm{kA}}\right)$}
\end{qoptions}
\raggedcolumns
\papersec{1}
\partsec{2}{CHEMISTRY}
\qn{19} Given are the information :\par
H$_2$(g) + $\dfrac12$S$_2$(g) = H$_2$S(g) ; K$^\circ_\mathrm P$ = 0.80\par
3H$_2$(g) + SO$_2$(g) = H$_2$S(g) + 2H$_2$O(g) ; K$^\circ_{\mathrm P_2}$ = 1.8$\times$10$^4$\par
The value of K$^\circ_\mathrm P$ for the reaction 4H$_2$(g) + 2SO$_2$(g) = S$_2$(g) + 4H$_2$O(g) is found to be.
\begin{qoptions}
\opt{A}{4.44$\times$10}
\opt{B}{2.25$\times$10$^4$}
\opt{C}{5.07$\times$10$^8$}
\opt{D}{4.26$\times$10$^6$}
\end{qoptions}
\qn{20} For the reaction NH$_3$(g) $\rightleftharpoons$ $\dfrac12$N$_2$(g) + $\dfrac32$H$_2$(g), the degree of dissociation in terms K$_\mathrm P$ is--
\begin{qoptions}
\opt{A}{$\alpha=\left(1+\dfrac{3\sqrt3}4\dfrac{\mathrm p}{\mathrm K_\mathrm P}\right)^{-\frac12}$}
\opt{B}{$\alpha=\left(1-\dfrac{3\sqrt3}4\dfrac{\mathrm p}{\mathrm K_\mathrm P}\right)^{-\frac12}$}
\opt{C}{$\alpha=\left(1+\dfrac{3\sqrt3}4\dfrac{\mathrm K_\mathrm P}{\mathrm P}\right)^{\frac12}$}
\opt{D}{$\alpha=\left(1-\dfrac{3\sqrt3}4\dfrac{\mathrm K_\mathrm P}{\mathrm P}\right)^{-\frac12}$}
\end{qoptions}
\qn{21} An aqueous solution is 0.3M in Al$^{3+}$ and 3M in Mg$^{2+}$.\par
Select the only \textbf{INCORRECT} option\par
Given : (K$_\mathrm{sp}$)$_{\mathrm{Al(OH)}_3}$ = 3 $\times$ 10$^{-34}$, (K$_\mathrm{sp}$)$_{\mathrm{Mg(OH)}_2}$ = 3 $\times$ 10$^{-12}$ M$^3$
\begin{qoptions}
\opt{A}{At pH = 9 both Al(OH)$_3$ and Mg(OH)$_2$ will precipitate}
\opt{B}{At pOH = 12 neither Al(OH)$_3$ nor Mg(OH)$_2$ will precipitate}
\opt{C}{For selective precipitation of Al(OH)$_3$ only, pH of solution should be 3 to 8.}
\opt{D}{Mg(OH)$_2$ will precipitate at maximum pH = 8}
\end{qoptions}
\qn{22} 1 mole each of A and D is introduced in 1 litre container. Simultaneously the following two equilibrium are established:\par
A $\rightleftharpoons$ B + C; K$_\mathrm C$ = 10$^6$ M\par
B + D $\rightleftharpoons$ A; K$_\mathrm C$ = 10$^{-6}$ M$^{-1}$\par
The equilibrium concentration of A will be
\begin{qoptions}
\opt{A}{10$^{-6}$ M}
\opt{B}{10$^{-3}$ M}
\opt{C}{10$^{-12}$ M}
\opt{D}{10$^{-4}$ M}
\end{qoptions}
\qn{23} 40 ml, 0.35 M NaOH solution is added to 50 ml 0.6 N H$_3$PO$_4$ solution. The pH of the mixture will be:\par
(Ka$_1$, Ka$_2$ and Ka$_3$ are 10$^{-3}$, 10$^{-8}$, 10$^{-12}$ respectively).
\begin{qoptions}
\opt{A}{7.82}
\opt{B}{5.76}
\opt{C}{8.84}
\opt{D}{9.32}
\end{qoptions}
\qn{24} For given two equilibriums attained in a container which are correct if degree of dissociation of A and A$'$ are $\alpha$ and $\alpha'$ respectively (P$'$ = Partial pressure at equilibrium)\par
A(s) $\rightleftharpoons$ 2B(g) + C(g); \quad K$_{\mathrm P_1}$ = 8 $\times$ 10$^{-2}$\par
A$'$(s) $\rightleftharpoons$ 2B(g) + D(g); \quad K$_{\mathrm P_2}$ = 2 $\times$ 10$^{-2}$
\begin{qoptions}
\opt{A}{$\dfrac{\mathrm K_{\mathrm P_2}}{\mathrm K_{\mathrm P_1}}=\left(\dfrac{3\alpha'+2\alpha}{3\alpha+2\alpha'}\right)^3\dfrac\alpha{\alpha'}$}
\opt{B}{$\dfrac{\mathrm P'_\mathrm C}{\mathrm P'_\mathrm D}$ = 4}
\opt{C}{P$'_\mathrm B$ = P$'_\mathrm C$ + P$'_\mathrm D$}
\opt{D}{$\alpha<\alpha'$}
\end{qoptions}
\qn{25} Which of the following statement are correct when equal volumes of 0.02 M HOCl \& 0.2 M CH$_3$COOH are mixed ?\par
K$_{\mathrm a_1}$ (HOCl) = 2$\times$10$^{-4}$; K$_{\mathrm a_2}$ (CH$_3$COOH) = 2$\times$10$^{-5}$ ; log 2 = 0.3
\begin{qoptions}
\opt{A}{$\alpha_1$ = 10$^{-1}$}
\opt{B}{$\alpha_2$ = 10$^{-2}$}
\opt{C}{pH = 2.7}
\opt{D}{[H$^+$] = 2 $\times$ 10$^{-3}$}
\end{qoptions}
\qn{26} For a reaction,\par
2A$_{(\mathrm g)}$ + 3B$_{(\mathrm g)}$ $\rightleftharpoons$ 4C$_{(\mathrm g)}$ + D$_{(\mathrm g)}$ \hfill K$^\circ_\mathrm C=\dfrac{32^2\times2}{27}$\par
If 4 moles of A are mixed with 6 moles of B in 2 litre container identify the correct options.
\begin{qoptions}
\opt{A}{At equilibrium, $\dfrac{[\mathrm C]}{[\mathrm B]}=\dfrac83$}
\opt{B}{At equilibrum, $\dfrac{[\mathrm D]}{[\mathrm A]}$ = 1}
\opt{C}{[A] \& [B] will be negligible as compared to [C] \& [D] at equilibrium.}
\opt{D}{K$_\mathrm P$ and K$_\mathrm C$ will be same for the above reaction.}
\end{qoptions}
\qn{27} Buffer capacity is defined in terms of buffer index which is the change in the concentration of buffer acid (or base) required for changing its pH value by one, keeping [C$_\mathrm{acid}$ + C$_\mathrm{salt}$] or [C$_\mathrm{base}$ + C$_\mathrm{salt}$] constant. The correct statements is/are for acidic buffer of HA and A$^-$
\begin{qoptions}
\opt{A}{Buffer capacity is highest when [HA] \& [A$^-$] are large and [HA] = [A$^-$]}
\opt{B}{Buffer capacity is highest when the pH is equal to (or very near) the pKa of the HA used.}
\opt{C}{For making a buffer solution with a pH = 2, among acetic acid [pKa = 4.74] \& chlorine acid [pKa = 1.95], acetic acid is better choice.}
\opt{D}{Among buffer A (0.1 M CH$_3$COOH \& 0.1 M CH$_3$COONa) and buffer B [0.2 M CH$_3$COOH \& 0.2 M CH$_3$COONa] buffer A is better buffer}
\end{qoptions}
\qn{28} For a gaseous reaction : A$_{(\mathrm g)}$ $\rightleftharpoons$ 3B$_{(\mathrm g)}$ + C$_{(\mathrm g)}$ , $\Delta$H is positive and the reaction attains equilibrium at 1 bar total pressure and 400 K. Identify the incorrect statements regarding the above reaction.
\begin{qoptions}
\opt{A}{On increase of temperature, equilibrium will be shifted in forward reaction.}
\opt{B}{When inert gas is introduced into a rigid container containing above equilibria equilibrium shifts towards left.}
\opt{C}{$\Delta$G$^\circ_{400}$ = 0 for the above reaction.}
\opt{D}{If Volume of vessel containing the above equilibria is increased without change in temperature then partial pressure of B decreases as compared to original equilibrium partial pressure of B.}
\end{qoptions}
\qn{29} A 1 gram pill contains morphine (M = 285.34), a monobasic compound whose conjugate acid has pKa = 8.21, in addition to an unknown amount of inert material. To analyze the morphine content of the pill, it is dissolved in 50 ml 0.1 M aq. HCl and the unreacted acid is back titrated with 0.1 M NaOH solution.\par
(I) If the pill contains only morphine, the titration will require 35 ml NaOH solution to reach the equivalence point. [(285.34)$^{-1}$ = 0.0035]\par
(II) Methyl red which changes colour from red to yellow between pH = 4.4 and pH = 6.2, would be a suitable indication for this titration, mentioned in statement(I).
\begin{qoptions}
\opt{A}{I is correct statement}
\opt{B}{I is incorrect statement}
\opt{C}{II is correct statement}
\opt{D}{II is incorrect}
\end{qoptions}
\qn{30} Which of the following statements are not correct ?
\begin{qoptions}
\opt{A}{The pH of pure water increases with increase of temperature.}
\opt{B}{The solubility of KAl(SO$_4$)$_2$ in terms of its solubility product is (K$_\mathrm{sp}$/2)$^{1/4}$.}
\opt{C}{The pH of a solution obtained by mixing 50 mL of 0.4 M HCl and 50 mL of 0.2 M NaOH would be one.}
\opt{D}{The standard dissociation constant of NH$_4$OH is 1.8 $\times$ 10$^{-5}$. The standard hydrolysis constant of NH$_4^+$ ions at 25$^\circ$C would be 5.56 $\times$ 10$^{-10}$.}
\end{qoptions}
\qn{31} For an equilibrium reaction, which of the following statements is/are correct ?
\begin{qoptions}
\opt{A}{If the reaction quotient of a reaction is greater that K$_\mathrm{eq}$, the reaction has a tendency to move in the backward direction.}
\opt{B}{If the reaction quotient of a reaction is greater than K$_\mathrm{eq}$, the reaction has a tendency to the move in the forward direction.}
\opt{C}{The addition of an inert gas at constant volume does not affect the extent of reaction.}
\opt{D}{The addition of an inert gas at constant pressure does not affect the extent of the reaction in general.}
\end{qoptions}
\qn{32} Following two equilibria are simultaneously established in a container\par
PCl$_5$(g) $\rightleftharpoons$ PCl$_3$(g) + Cl$_2$ ; CO(g) + Cl$_2$(g) $\rightleftharpoons$ COCl$_2$(g)\par
If some Ni(s) is introduced in the container establishing another equilibrium.\par
Ni(s) + 4CO(g) $\rightleftharpoons$ [Ni(CO)$_4$](g)\par
Which among the following statements is/are definitely true about this sytem ?
\begin{qoptions}
\opt{A}{COCl$_2$ concentration decreases as the new equilibrium is established.}
\opt{B}{As Ni is a solid it will not have any impact on the existing two equilibrium reactions.}
\opt{C}{PCl$_3$ concentration decreases to re--attain the equilibrium.}
\opt{D}{PCl$_5$ concentration increases to attain the new equilibrium.}
\end{qoptions}
\paraq{33}{34}Titration of diprotic acid (H$_2$A) by strong base has been summarised in following graph\par
\qfig{p1_chem_s0_q32_1}
\qn{33} What are the pH of solution when it shows maximum buffer action ?
\begin{qoptions}
\opt{A}{2.34, 6}
\opt{B}{6, 9.66}
\opt{C}{2.34, 9.66}
\opt{D}{6, 12}
\end{qoptions}
\qn{34} What is the pH at II equivalence point?
\begin{qoptions}
\opt{A}{6.0}
\opt{B}{11.83}
\opt{C}{9.66}
\opt{D}{12}
\end{qoptions}
\paraq{35}{36}In a closed rigid vessel, N$_2$ and H$_2$ gases are taken in 9 : 13 mole ratio, by which the following equilibria are estab--lished :\par
N$_2$(g) + 3H$_2$(g) $\rightleftharpoons$ 2NH$_3$(g);K$_{\mathrm P_1}$ = ?\par
N$_2$(g) + 2H$_2$(g) $\rightleftharpoons$ N$_2$H$_4$(g);K$_{\mathrm P_2}$ = ?\par
At equilibrium, the total pressure is 14 atm, the partial pressure of ammonia is 2 atm and the partial pressure of hydrogen is 4 atm.
\qn{35} The value of K$_{\mathrm P_1}$ is :
\begin{qoptions}
\opt{A}{$\dfrac1{80}$atm$^2$}
\opt{B}{$\dfrac1{80}$atm$^{-2}$}
\opt{C}{80 atm$^2$}
\opt{D}{$\dfrac1{40}$atm$^{-2}$}
\end{qoptions}
\qn{36} The value of $\dfrac1{\mathrm K_{\mathrm P_2}}$ is :
\begin{qoptions}
\opt{A}{$\dfrac3{80}$atm$^{-2}$}
\opt{B}{$\dfrac{80}3$atm$^2$}
\opt{C}{$\dfrac{40}3$atm$^2$}
\opt{D}{$\dfrac53$atm$^{-2}$}
\end{qoptions}
\raggedcolumns
\papersec{1}
\partsec{3}{MATHEMATICS}
\qn{37} On Argand plane, let A be the total number of integral points inside the region $|z|<4$ and B the total number of integral points satisfying $|z-3i|<4$ and $\dfrac\pi6<\arg(z-3i)<\dfrac{5\pi}6$ then the value of $\dfrac AB$ is :
\begin{qoptions}
\opt{A}{2}
\opt{B}{1}
\opt{C}{5}
\opt{D}{None}
\end{qoptions}
\qn{38} If A($z_1$); B($z_2$) and C($z_3$) are vertices of a triangle such that $z_3=\dfrac{z_2-iz_1}{1-i}$ and $|z_1|=3$; $|z_2|=4$ and $|z_2+iz_1|=|z_1|+|z_2|$ then area of $\Delta$ABC is :
\begin{qoptions}
\opt{A}{6}
\opt{B}{12}
\opt{C}{5}
\opt{D}{6.25}
\end{qoptions}
\qn{39} If $\alpha$ is an imaginary fifth root of unity then the value of
$\log_{\left(\frac13\right)^{1/3}}\!\left[\log_{\frac18}\left|\dfrac1{1+\alpha+\alpha^2+\alpha^3-\dfrac1\alpha}\right|\right]$ is :
\begin{qoptions}
\opt{A}{2}
\opt{B}{3}
\opt{C}{4}
\opt{D}{9}
\end{qoptions}
\qn{40} $\displaystyle\sum_{k=0}^{7}\left(\dfrac{{}^7\mathrm C_k}{{}^{14}\mathrm C_k}\sum_{r=k}^{14}\left({}^r\mathrm C_k\right)\left({}^{14}\mathrm C_r\right)\right)=(k)^{k+1}$ then $k=$ ?
\begin{qoptions}
\opt{A}{14}
\opt{B}{4}
\opt{C}{6}
\opt{D}{7}
\end{qoptions}
\qn{41} Let $f(n)=\cot^2\dfrac\pi n+\cot^2\dfrac{2\pi}n+\dots\dots+\cot^2\dfrac{(n-1)\pi}n$ (where $n>1$, $n\in\mathrm N$, $n$ odd number), then $\displaystyle\lim_{n\to\infty}\dfrac{3f(n)}{(n+1)(n+2)}$ is
\begin{qoptions}
\opt{A}{$\dfrac12$}
\opt{B}{$\dfrac34$}
\opt{C}{2}
\opt{D}{1}
\end{qoptions}
\qn{42} Let $z_1,z_2,z_3$ be complex numbers such that $|z_1|=|z_2|=|z_3|=10$ and $z_2\ne z_3$ then $\dfrac{\displaystyle\min_{a\in\mathrm R}\left|az_2+(1-a)z_3-z_1\right|}{|z_1-z_2||z_1-z_3|}=$ :
\begin{qoptions}
\opt{A}{0.05}
\opt{B}{0.5}
\opt{C}{0.025}
\opt{D}{0.25}
\end{qoptions}
\qn{43} If $a=1+\dfrac{\omega^2x^2}{2!}+\dfrac{\omega^4x^4}{4!}+\dots..$ \& $b=-x-\dfrac{\omega^2x^3}{3!}-\dfrac{\omega^4x^5}{5!}\dots..$ where $\omega$ is non real cube root of unity then which of the following is incorrect ?
\begin{qoptions}
\opt{A}{$a^2-b^2\omega^2=1$}
\opt{B}{$a=\dfrac{e^{\omega x}-e^{-\omega x}}2$}
\opt{C}{$a^2+b^2\omega^2+1=0$}
\opt{D}{$(a^{1/3}+b^{1/3})^{2/3}=1$}
\end{qoptions}
\qn{44} Let $z_1,z_2,z_3$ be the coordinates of vertices A, B, C of the triangle if $\omega_1=z_1-z_2$; $\omega_2=z_3-z_2$ and $\angle$A $=90^\circ$ then :
\begin{qoptions}
\opt{A}{$\mathrm{Re}\left(\dfrac{\omega_2}{\omega_1}\right)=1$}
\opt{B}{$\mathrm{Re}(\omega_1\overline{\omega_2})=|\omega_1|^2$}
\opt{C}{$\mathrm{Re}(\omega_1\overline{\omega_2})=|\overline{\omega_2}|^2$}
\opt{D}{$\mathrm{Re}\left(\dfrac{\omega_1}{\omega_2}\right)=1$}
\end{qoptions}
\qn{45} If curve A $\equiv\{z:|z-3|+|z+3|=8\}$; and B $\equiv\{z:|z-3|=k\}$, $k\in\mathrm R^+$ touches the curve `A' internally; and given another curve C $\equiv\{z:\|z-3|-|z+3\|=4\}$ then :
\begin{qoptions}
\opt{A}{A and C intersect orthogonally and $k=1$}
\opt{B}{$k=1$}
\opt{C}{$k=\dfrac32$}
\opt{D}{B also touches C internally}
\end{qoptions}
\qn{46} The coefficient of $x^{20}$ in $\displaystyle\sum_{m=0}^{18}{}^{18}\mathrm C_m(x-2)^{18-m}3^m+\sum_{m=0}^{19}{}^{19}\mathrm C_m(x-2)^{19-m}3^m+\sum_{m=0}^{20}{}^{20}\mathrm C_m(x-2)^{20-m}3^m+\dots\dots+\sum_{m=0}^{30}{}^{30}\mathrm C_m(x-2)^{30-m}3^m$ cannot be :
\begin{qoptions}
\opt{A}{${}^{31}\mathrm C_{20}$}
\opt{B}{${}^{31}\mathrm C_{21}$}
\opt{C}{${}^{31}\mathrm C_{19}$}
\opt{D}{${}^{30}\mathrm C_{19}$}
\end{qoptions}
\qn{47} $\displaystyle\sum_{r=0}^{99}\left(\dfrac{101}{100}\right)\left(\dfrac{r\,{}^{100}\mathrm C_r\,{}^{100}\mathrm C_{r+1}}{r+2}\right)={}^x\mathrm C_y$ then $x+y$ can be :
\begin{qoptions}
\opt{A}{298}
\opt{B}{302}
\opt{C}{114}
\opt{D}{162}
\end{qoptions}
\qn{48} $(x+m)^m-{}^m\mathrm C_1(x+m-1)^m+{}^m\mathrm C_2(x+m-2)^m\dots\dots+(-1)^mx^m$ when $m=6$ equal to `$\ell$' then `$\ell$' is divisible by:
\begin{qoptions}
\opt{A}{7}
\opt{B}{8}
\opt{C}{9}
\opt{D}{11}
\end{qoptions}
\qn{49} $\mathrm N=\displaystyle\sum_{a=0}^{50}\sum_{b=0}^{a}\sum_{c=0}^{a}\sum_{d=0}^{50-a}\sum_{e=0}^{c+d}\sum_{f=0}^{50-a-d}$\par$\dbinom{50}a\dbinom ab\dbinom ac\dbinom{50-a}d\dbinom{c+d}e\dbinom{50-a-d}f$\par
Which of the following are correct :
\begin{qoptions}
\opt{A}{Units digit of N is 0}
\opt{B}{Number of digit in N are 51}
\opt{C}{Total number of divisors of N are 2801}
\opt{D}{Total number of zeroes in number N are 50}
\end{qoptions}
\qn{50} Let $\mathrm S_n=1-\dfrac12+\dfrac13-\dfrac14+\dots..+(-1)^{n-1}\dfrac1n$ then :
\begin{qoptions}
\opt{A}{$\displaystyle\sum_{r=1}^{n}\dfrac1{n+r}=\mathrm S(2n)$}
\opt{B}{$\mathrm S(2n)<1\ \forall\ n\in\mathrm N$}
\opt{C}{$\mathrm S(2n)\ge0.5\ \forall\ n\in\mathrm N$}
\opt{D}{$0.5<\mathrm S(n)<1\ \forall\ n\in\mathrm N$}
\end{qoptions}
\paraq{51}{52}Let the roots of $x^j=1$ be, $\omega_{j1},\omega_{j2},\omega_{j3}\dots\dots\omega_{j(j-1)}$ and $(1-\omega_{j1})(1-\omega_{j2})\dots\dots(1-\omega_{j(j-1)})=j$
\qn{51} Find the value of $\displaystyle\sum_{j=2}^{61}\sum_{i=1}^{j-1}\dfrac1{1-\omega_{ji}}$ :
\begin{qoptions}
\opt{A}{915}
\opt{B}{$-915$}
\opt{C}{1830}
\opt{D}{None of these}
\end{qoptions}
\qn{52} The value of $\displaystyle\sum_{r=1}^{j-1}r(\omega_{jr}+\omega_{j(j-r)})$ equals :
\begin{qoptions}
\opt{A}{$-j$}
\opt{B}{$j$}
\opt{C}{$-j-1$}
\opt{D}{$-\dfrac j2$}
\end{qoptions}
\paraq{53}{54}If $(1+px+x^2)^n=1+a_1x+a_2x^2+\dots..a_{2n}x^{2n}$; $p>3$, $p\in\mathrm N$
\qn{53} Which of the following is true for $1<r<2n$
\begin{qoptions}
\opt{A}{$(np+pr)a_r=(r+1)a_{r+1}+(r-1)a_{r+1}$}
\opt{B}{$(np-pr)a_r=(r+1)a_{r+1}+(r-1-2n)a_{r-1}$}
\opt{C}{$(np-pr)a_r=(r+1)a_{r+1}+(r-1-n)a_{r-1}$}
\opt{D}{$(2np-pr)a_r=(r+1+n)a_{r+1}+(r+1-n)a_{r-1}$}
\end{qoptions}
\qn{54} The remainder obtained when $a_1+5a_2+9a_3+13a_4+\dots..+(8n-3)a_{2n}$ is divided by $(p+2)$ is :
\begin{qoptions}
\opt{A}{1}
\opt{B}{2}
\opt{C}{3}
\opt{D}{0}
\end{qoptions}
\end{multicols}
\end{document}
